% Root/B/C/D full actual-file audits passed; auxiliary unbounded-period input.
\begin{proposition}[Proved: Brauer residues and period-degree covers for one exceptional curve]
\label{D:single-curve-Brauer-residue-cover}
Let $(R,\mathfrak m,k)$ be a complete excellent normal local
$k$-algebra of dimension two, with algebraically closed
coefficient and residue field $k$ of characteristic $p>0$.
Put $U=\operatorname{Spec}R-\{\mathfrak m\}$.  Suppose that
there is a projective regular resolution
$\pi:Y\to\operatorname{Spec}R$, an isomorphism over $U$,
whose reduced exceptional fibre is one smooth connected
curve $C$.  Write $g=g(C)$ and $C^2=-d<0$.
For every positive integer $n$ prime to $p$, residue gives
a canonical isomorphism
\[
       \operatorname{Br}'(U)[n]
          \simeq H^1(C,\mathbf Z/n\mathbf Z),
       \qquad |\operatorname{Br}'(U)[n]|=n^{2g},       \tag{D.RC.1}
\]
where $\operatorname{Br}'$ denotes cohomological torsion
Brauer classes.

If $\alpha$ has order $m$ prime to $p$, there is a connected
cyclic finite etale cover $Y'\to Y$ of degree exactly $m$
such that its restriction $V\to U$ kills $\alpha$.
The cover $V$ is the puncture of a complete excellent normal
local ring $R'$ finite over $R$.  The induced resolution
$Y'\to\operatorname{Spec}R'$ has a single smooth connected
exceptional curve $C'$, with
\[
          g(C')-1=m(g-1),\qquad (C')^2=-md.          \tag{D.RC.2}
\]
Thus the splitting cover is etale on the entire resolution,
not only on the puncture.  No assertion about a representation
count uniform in $m$ is part of this proposition.
\end{proposition}

\begin{proof}
The Picard lifting argument in
Lemma~\ref{D:local-Picard-prime-to-p-generators} gives a
surjection $\operatorname{Pic}(Y)\to\operatorname{Pic}(C)$:
the obstructions through the exceptional thickenings lie
in coherent $H^2$ on a curve, and Grothendieck existence
algebraizes a compatible line bundle and its inverse.
Since $R$ is strictly henselian, proper base change gives
\[
           H^a(Y,\mu_n)\simeq H^a(C,\mu_n).
                                                        \tag{D.RC.3}
\]
This is precisely the scope checked in that lemma, including
a possibly nonreduced scheme-theoretic closed fibre.
For a smooth proper curve over $k$, its degree map identifies
$\operatorname{Pic}(C)/n$ with $H^2(C,\mu_n)=\mathbf Z/n$,
and $H^3(C,\mu_n)=0$
\citep[Lemma~59.69.1 and complete proof, Tag 03RQ]
{StacksSingleCurveBrauerD}.
The Kummer square for restriction to $C$ therefore shows
\[
       H^2(Y,\mu_n)=\operatorname{Pic}(Y)/n,
       \qquad \operatorname{Br}'(Y)[n]=0,
       \qquad H^3(Y,\mu_n)=0.                         \tag{D.RC.4}
\]

Also $\operatorname{Pic}(Y)\to\operatorname{Pic}(U)$ is
surjective.  Indeed every line bundle on the regular integral
scheme $U$ is represented by a Cartier divisor; its Weil
divisor closure on the regular scheme $Y$ is Cartier.
Consequently the image of $H^2(Y,\mu_n)$ in $H^2(U,\mu_n)$
is exactly the Kummer image of $\operatorname{Pic}(U)/n$.
The localization sequence and (D.RC.4) give
\[
 \operatorname{Br}'(U)[n]
  =H^2(U,\mu_n)/(\operatorname{Pic}(U)/n)
  \simeq H^3_C(Y,\mu_n)
  \simeq H^1(C,\mathbf Z/n\mathbf Z).                 \tag{D.RC.5}
\]
The last isomorphism is absolute purity for the regular
codimension-one immersion $C\hookrightarrow Y$; its Tate
twist cancels the twist in $\mu_n$.
We use the exact regular-scheme scope of
\citep[Expose~XVI, Theorem~3.1.1]{A_ILO2014}, as in
Theorem~\ref{D:local-surface-constant-H2}; no finite-type
over $k$ hypothesis is imposed on $Y$.
The last group has order $n^{2g}$ by the same curve
cohomology calculation.  The construction is functorial
under finite etale base change of the pair $(Y,C)$.

Now take $n=m$ and let $\chi\in H^1(C,\mathbf Z/m)$
be the residue of $\alpha$.  The isomorphism (D.RC.5)
preserves orders, so the corresponding homomorphism
$\pi_1(C)\to\mathbf Z/m$ is surjective.  Its torsor
$C'\to C$ is therefore connected and has degree $m$.
The proper henselian equivalence for finite etale covers
\citep[Lemma~58.9.1 and complete proof, Tag 0A48]
{StacksSingleCurveBrauerD}
lifts this torsor, including its cyclic action, to a finite
etale $Y'\to Y$.  This cover is connected: any nonempty
open and closed component is proper over $R$ and hence
meets the closed fibre, whereas $C'$ is connected.
It is regular, and thus integral.

Let $L=k(Y')$, and normalize $R$ in $L$ to obtain $R'$.
This is finite by excellence, complete, and local because
it is a finite domain over a henselian local ring.  Its
residue field is $k$.  Equivalently it is the Stein ring
$H^0(Y',\mathcal O_{Y'})$: that ring is finite over $R$,
normal, and has fraction field $L$, so it is exactly the
integral closure.  The map $Y'\to\operatorname{Spec}R'$
is proper and birational, and is an isomorphism away from
the closed point, where it agrees with the finite normal
cover $V\to U$.  Its reduced closed fibre is $C'$.

Pullback of (D.RC.5) along $Y'\to Y$ is the corresponding
residue isomorphism for this new pair.  The pulled-back
character $\chi|_{C'}$ is zero by the defining torsor, so
$\alpha|_V=0$.  Finally $C'\to C$ is etale of degree $m$.
Riemann--Hurwitz gives the first identity in (D.RC.2)
\citep[Lemma~53.12.2, Tag 0C1D]{StacksSingleCurveBrauerD}.
The normal line bundle of $C'$ in $Y'$ is the pullback of
the normal line bundle of $C$ in $Y$.  Its degree is $mC^2$,
which gives the second identity.
\end{proof}

\begin{remark}[Quantitative issue still to be proved]
The splitting degree equals the geometric period, but it
can grow.  Applying a rank-$r$ count on the varying cover
without controlling its constants does not give the desired
uniform rank-quadratic bound.  A residue character of order
$m$ suggests rank-$r/m$ inertia blocks over $C'$, but normal
coordinate powers shift their weights.  A description on
the reduced curve alone supplies no formal-neighbourhood
decomposition or Artinian parameter presentation.
\end{remark}
